\section{Introducci\'on}
En el día a día estamos rodeados por sistemas dinámicos que
evolucionan en el tiempo.
Cada uno de estos sistemas puede ser representado matemáticamente
mediante lo que suele llamarse un «modelo».
El modelo no es más que una simplificación de un sistema real,
tomando los posibles elementos que interactúan y representar
matemáticamente dicha interacción para predecir el estado de dicho
sistema en un momento del tiempo.
Los sistemas cuyo modelado es más simple, son los sistemas LTI
o lineales invariantes en el tiempo.
A partir de estos modelos se ha desarrollado una sólida teoría sobre
el análisis, observación y control de sistemas dinámicos que ha beneficiado
en diferentes áreas, tanto del conocimiento como de la industria.

